\chapter{Casos de Estudio Reales}
\label{sec.cap8:casos_estudio}

\section{Caso 1: Depósito Corporativo Iron Mountain (Barracas, Buenos Aires, 2014)}
\label{sec.cap8:iron_mountain}

El 5 de febrero de 2014 se originó un incendio catastrófico en la planta de resguardo y archivo de documentación bancaria
y corporativa de la empresa Iron Mountain, ubicada en el barrio de Barracas, Ciudad Autónoma de Buenos Aires.

\begin{itemize}
    \item Carga de Fuego y Configuración Espacial: La instalación albergaba una carga de fuego extraordinaria ($Q_f \gg
          150\,\mathrm{kg/m^2}$), estructurada en racks metálicos autoportantes de más de diez metros de altura densamente
          cargados con cajas de cartón corrugado y papel comprimido. Esta disposición vertical generó chimeneas de tiro
          térmico ascendente que aceleraron la velocidad de propagación de las llamas hacia la cubierta del edificio.
    \item Falla de Barreras Activas y Pasivas: Las pericias judiciales determinaron que el sistema de rociadores
          automáticos (diseñado bajo lineamientos NFPA 13) se encontraba inoperativo al momento del inicio del fuego: las
          válvulas sectoriales estaban cerradas o carecían de presión hídrica de servicio, y las bombas principales no
          arrancaron ante la demanda. Simultáneamente, la estructura portante de acero carecía de revestimientos ignífugos
          pasivos proyectados (pinturas intumescentes o morteros de vermiculita) que garantizaran una resistencia $F120$ o
          $F180$ acorde a la carga térmica presente.
    \item Colapso Estructural y Víctimas: Al alcanzarse temperaturas en cabriadas superiores a los $900\Celsius$, el
          acero experimentó fluencia térmica y pérdida de su límite elástico. El pandeo de los perfiles del techo ejerció
          fuerzas horizontales de tracción y empuje asimétrico sobre los muros perimetrales de ladrillo macizo. La pared
          este colapsó íntegramente hacia la calle Jovellanos, sepultando a una dotación de diez rescatistas (ocho bomberos
          de la Policía Federal Argentina y dos agentes de Defensa Civil).
    \item Lecciones Aprendidas: El siniestro evidenció que la instalación nominal de sistemas contra incendio carece de
          validez sin auditorías periódicas, protocolos de inspección inopinada y circuitos cerrados de presurización
          supervisados por paneles centrales FACP. Asimismo, demostró la necesidad de delimitar zonas de seguridad
          perimetral (distancia de derrumbe igual a $1.5$ veces la altura del muro) ante incendios en estructuras metálicas
          desnudas.
\end{itemize}

\section{Caso 2: Subestación Transformadora Caballito de Edesur (Buenos Aires, 2024)}
\label{sec.cap8:edesur_caballito}

El 10 de febrero de 2024 se desató un voraz incendio en la subestación Caballito de la compañía distribuidora de energía
Edesur, dejando sin suministro eléctrico a más de $60000$ usuarios residenciales e industriales y afectando el bombeo de
agua de la empresa AySA en la Ciudad Autónoma de Buenos Aires.

\begin{itemize}
    \item Origen de la Falla y Riesgo Asociado: El siniestro se inició durante maniobras de mantenimiento sobre una
          máquina portátil de filtrado y secado al vacío de aceite aislante conectada a un transformador de potencia de
          $132/13.2\,\mathrm{kV}$. La rotura de un conducto flexible bajo presión pulverizó aceite mineral caliente sobre
          partes energizadas o superficies a alta temperatura, desatando una combustión violenta de combustible líquido
          (fuego Clase B) en estrecha proximidad con equipamiento de alta tensión (riesgo Clase C).
    \item Propagación por Canalizaciones: Las llamas se extendieron con rapidez a través de las canalizaciones subterráneas
          de conductores de media y baja tensión. La presencia de aislaciones convencionales de PVC no ignífugas produjo una
          densa humareda negra cargada de hollín y gas clorhídrico, dificultando el acceso visual y operativo de las
          fuerzas de rescate.
    \item Estrategia de Extinción Aplicada: Doce dotaciones de Bomberos de la Ciudad combatieron el siniestro durante más
          de seis horas. Debido a la naturaleza líquida del hidrocarburo en combustión, el agua directa resultaba inviable
          por el peligro de dispersión del combustible por flotación. Se desplegaron líneas de ataque con espuma mecánica
          AFFF (formadora de película acuosa) para sellar la superficie del aceite, sofocar la masa ígnea e inhibir la
          reignición, complementadas con cortinas de agua pulverizada para refrigerar las cubas de los transformadores
          adyacentes y evitar su rotura por sobrepresión interna.
    \item Lecciones para la Ingeniería: El evento remarcó la importancia de construir muros cortafuego $F120$ entre
          transformadores de potencia contiguos, instalar piletas colectoras de derrame con lechos de grava apagallamas y
          reemplazar conductores antiguos por cables con aislación libre de halógenos (IRAM 62267) según lo estipulado por
          la reglamentación AEA 90364.
\end{itemize}

\section{Caso 3: Fábrica de Semiconductores Renesas Electronics (Naka, Japón, 2021)}
\label{sec.cap8:renesas_cleanroom}

El 19 de marzo de 2021 se produjo un incendio en la sala limpia (cleanroom de $300\,\mathrm{mm}$) de la planta fabril de
semiconductores de Renesas Electronics en Naka, Japón, uno de los principales proveedores globales de microcontroladores
para la industria automotriz.

\begin{itemize}
    \item Causa Raíz Eléctrica: Una sobrecorriente prolongada en la línea de suministro de un equipo galvánico de
          galvanoplastia y electrodeposición de cobre generó disipación por efecto Joule en las conexiones de potencia. El
          calentamiento localizado no fue despejado oportunamente por las protecciones del tablero secundario, alcanzando
          la temperatura de ignición de las resinas poliméricas que conformaban la envolvente del equipo.
    \item Dispersión del Humo y Daño a Activos Electrónicos: Aunque la extensión del fuego directo quedó restringida a
          un área de $600\,\mathrm{m^2}$, el sistema de climatización y presurización de la sala limpia (destinado a mantener
          el ambiente libre de polvo con filtros HEPA) succionó los subproductos pirolíticos y los recirculó a través de
          toda la nave de producción. Las partículas submicrónicas de hollín, combinadas con gases clorados y sulfurosos
          desprendidos de los plásticos, se depositaron sobre las lentes ópticas de los equipos de litografía ultravioleta y
          las cámaras de deposición química en fase de vapor (CVD).
    \item Impacto Económico Global: La corrosión y contaminación superficial forzó el desmantelamiento de veintitrés
          máquinas de producción crítica y demoró tres meses la reactivación fabril, agravando la escasez mundial de chips
          en la industria automotriz con pérdidas valoradas en cientos de millones de dólares.
    \item Implicancias Tecnológicas: Este desastre constituye el ejemplo paradigmático de la vulnerabilidad de las
          instalaciones electrónicas modernas frente al humo. En tales recintos, el combate con agua resulta inadmisible por
          ruina de los equipos. Resulta imperativo implementar detección por aspiración de alta sensibilidad (ASD/VESDA)
          interconectada a los ductos de retorno del sistema HVAC para corte automático de ventilación, asociada a sistemas
          fijos de inundación total con agentes limpios (FK-5-1-12 o Inergen bajo norma NFPA 2001) y relés de arco AFDD.
\end{itemize}

